\section{Model Architecture}
\label{supp:sec:model_architecture}

In this section, we provide comprehensive details about the Transformer model architectures considered in this work.
We implement all models in PyTorch~\citep{Paszke2019PyTorch} and adapt the implementation of Transformer-XL from VPT~\citep{openai2022vpt}.

\subsection{Observation Encoding}
Experiments conducted on both DMLab and RoboMimic include RGB image observations.
For models trained on DMLab, we use a ConvNet~\citep{he2015resnet} similar to the one used in \citet{espeholt2018impala}. 
For models trained on RoboMimic, we follow \citet{mandlekar2021matters} to use a ResNet-18 network~\citep{he2015resnet} followed by a spatial-softmax layer~\citep{finn2015deep}. We use independent and separate encoders for images taken from the wrist camera and frontal camera.
Detailed model parameters are listed in Table~\ref{supp:table:vision_encoder}.

\begin{table}[h]
\caption{Model hyperparameters for vision encoders.}
\vspace{0.1in}
\centering
\begin{tabular}{@{}ll@{}}
\toprule
\textbf{Hyperparameter} & \textbf{Value}    \\ \midrule
\multicolumn{2}{c}{DMLab} \\\midrule
Image Size     & 72 $\times$ 96 \\
Number of ConvNet Blocks     & 1       \\
Channels per Block      & [16, 32, 32]      \\
Output Size     & 256        \\\midrule
\multicolumn{2}{c}{RoboMimic} \\\midrule
Image Size     & 84 $\times$ 84 \\
Random Crop Height     & 76       \\
Random Crop Width     & 76       \\
Number of Randomly Cropped Patches     & 1       \\
ConvNet Backbone      & ResNet-18~\citep{he2015resnet}      \\
Output Size     & 64        \\
Spatial-Softmax Number of Keypoints     & 32        \\
Spatial-Softmax Temperature     & 1.0        \\
Output Size     & 64        \\\bottomrule
\end{tabular}
\label{supp:table:vision_encoder}
\end{table}


Since DMLab is highly partially observable, we follow previous work~\citep{espeholt2018impala,fan2022minedojo,openai2022vpt} to supply the model with previous action input.
We learn 16-dim embedding vectors for all discrete actions.

To encode proprioceptive measurement in RoboMimic, we follow \citet{mandlekar2021matters} to not apply any learned encoding. Instead, these types of observation are concatenated with image features and passed altogether to the following layers.
Note that we do not provide previous action inputs in RoboMimic, since we find doing so would incur significant overfitting.

\subsection{Transformer Backbone}
We use Transformer-XL~\citep{dai2019transformerxl} as our model backbone, adapted from \citet{openai2022vpt}. Transformer-XL splits long sequences into shorter sub-sequences that reduce the computational cost of attention while allowing the hidden states to be carried across the entire input by attending to previous keys and values.
This feature is critical for the long sequence inputs necessary for cross-episodic attention.
Detailed model parameters are listed in Table~\ref{supp:table:xf_xl}.

\begin{table}[ht]
\caption{Model hyperparameters for Transformer-XL.}
\vspace{0.1in}
\centering
\begin{tabular}{@{}lll@{}}
\toprule
\textbf{Hyperparameter} & \textbf{Value (DMLab)} & \textbf{Value (RoboMimic)} \\ \midrule
Hidden Size             & 256                    & 400                        \\
Number of Layers        & 4                      & 2                          \\
Number of Heads         & 8                      & 8                          \\
Pointwise Ratio         & 4                      & 4                          \\ \bottomrule
\end{tabular}
\label{supp:table:xf_xl}
\end{table}

\subsection{Action Decoding}
To decode joystick actions in DMLab tasks, we learn a 3-layer MLP whose output directly parameterizes a categorical distribution. This action head has a hidden dimension of 128 with ReLU activations. 
The ``Goal Maze'' and ``Irreversible Path'' tasks have an action dimension of 7, while ``Watermaze'' has 15 actions.
To decode continuous actions in RoboMimic, we learn a 2-layer MLP that parameterizes a Gaussian Mixture Model (GMM) with $5$ modes that generates a 7-dimensional action.
This network has a hidden dimension of 400 with ReLU activations.
During deployment, we employ the ``low-noise evaluation'' trick~\citep{hoffman2020acme}.
